%% kapa: exam=1 round=37 year=2026 period=2 form=A subject=appraisal_law subject_ko=감정평가관계법규 seq=5 number=5
%% source: Q-Net 감정평가사 기출문제 (cst00309) artlSeq=5258924 file=raw/1008/5258924/2026년 제37회 감정평가사 1차 2교시.pdf
%% conversion: kordoc 4.24.51 -> tools/gen1.py (auto); review=proofread-sonnet; flags=none
\begin{question}{5}
국토의 계획 및 이용에 관한 법령상 용도지역 및 용도구역에서의 행위 제한에 관한 설명으로 옳지 않은 것은?
\begin{choices}
  \item 도시지역에 대하여는 ｢도로법｣ 제40조에 따른 접도구역 규정을 적용한다.
  \item 공업지역에 대하여는 ｢농지법｣ 제8조에 따른 농지취득자격증명 규정을 적용하지 아니한다.
  \item 도시혁신구역에 대하여는 ｢문화예술진흥법｣ 제9조에 따른 건축물에 대한 미술작품의 설치를 도시혁신계획으로 따로 정할 수 있다.
  \item 복합용도구역으로 지정된 지역은 ｢건축법｣ 제69조에 따른 특별건축구역으로 지정된 것으로 본다.
  \item 도시혁신구역에 대하여는 ｢주차장법｣ 제19조에 따른 부설주차장의 설치를 도시혁신계획으로 따로 정할 수 있다.
\end{choices}
\end{question}
\answer{1}
